\section{Introduction}

Cislunar space is the region in and around the Earth--Moon system. A
trajectory-correction maneuver, or correction burn, deliberately changes a
spacecraft's velocity to keep its path within mission constraints. The required
change is summarized by delta-v (measured in m/s). Delta-v is not propellant
mass, although it is a standard measure of maneuver effort and is related to
fuel use through the vehicle and propulsion model.

For a crew-carrying mission, correction planning is a coupled decision problem:
burn timing and delta-v interact with navigation uncertainty, maneuver
execution, crew timelines, entry conditions, and reserve constraints. Public
Artemis trajectory-design research therefore treats robust correction cost,
meaning correction delta-v evaluated with modeled uncertainty, together with
mission constraints rather than as an unconstrained curve-fitting task
\cite{woffinden2023}. A learning \emph{surrogate}, a fast predictive model that
approximates a more expensive simulation, could help screen candidate scenarios.
That possibility is a motivation for testing, not evidence that a particular
machine-learning model is useful or safe.

The paper adopts a human-centered decision-support framing. A model score is
not the decision: a responsible reviewer needs to know which source data,
physical assumptions, uncertainty boundary, and failure cases produced the
score before deciding whether a model warrants further study. Crew-related
requirements therefore remain separate from predictive performance, consistent
with the intended-use and human-system boundaries in the applicable U.S.
National Aeronautics and Space Administration (NASA) standards
\cite{nasa7009,nasa3001}. This work does not conduct a user study or measure
usability, trust, workload, or operator performance. From a human--computer
interaction (HCI) perspective, its contribution is the evidence design around
a possible future decision-support tool.

Quantum machine learning (QML) uses quantum-state representations or quantum
circuits inside a learning method. In a quantum-kernel model, a circuit maps
each classical input to a quantum state and the overlap between two states is
used as a similarity score. This offers feature maps different from common
classical models \cite{schuld2019,havlicek2019}, but difference alone does not
imply usefulness. Encoding data introduces loading choices, similarity
geometry, and resource costs. A circuit may be difficult to simulate in general
while still being approximated well by a classical learner on the tested data
\cite{huang2021,sweke2025}. Conversely, a favorable result can arise from a
weakly tuned comparator, related mission cases leaking across a data split,
thresholds chosen after seeing outcomes, or omitted failures. These are
benchmarking problems before they are quantum-computing problems.

This paper documents a complete public-data experimental program for
trajectory-correction prediction and feasibility screening. The purpose was
not to assert a priori that QML would win. Its purpose was to establish a
scientifically usable decision process: define a bounded task, verify the
simulator against eligible public endpoints, freeze the development protocol,
compare QML with strong controls, record negative outcomes, and prohibit
claims that the evidence cannot support. This stance is consistent with model
and simulation credibility practice \cite{nasa7009} and with recent calls for
fair QML comparisons \cite{bowles2024,schnabel2024}.

The central research questions were deliberately narrow:
\begin{enumerate}
\item Can the primary fidelity-style quantum kernel (repository identifier
Q01) improve correction-cost prediction over physics-derived and classical
controls under the frozen development protocol?
\item If Q01 fails, can the projected follow-up Q01b improve the same task
without changing data access or data splits?
\item Can the feasibility-only quantum kernel (FQK) meet its fixed screening
criteria for identifying independently propagated feasible cases?
\item Which conclusions remain valid when calibration, final-test,
mission-loop, and hardware evidence stay locked?
\end{enumerate}

The answer to the first three questions is negative for the tested models and
settings. The fourth question is the paper's main methodological contribution.
The work establishes a transparent evidence hierarchy that distinguishes
simulator credibility, development-only model evidence, future-design lessons,
and prohibited operational claims. Negative results are not treated as failed
record keeping; they are the result when a hypothesis does not meet its
predeclared standard.

The paper makes six contributions:
\begin{enumerate}
\item A source-bound cislunar benchmark protocol, meaning that inputs are tied
to identified public files and recorded source versions, with grouped
development folds, training-fold-only transformations, and retained numerical
feasibility outcomes.
\item A simulator credibility chain that reports its eligible scope and its
source-ineligible exception rather than converting missing evidence into a
pass or failure.
\item A preregistered QML comparison against physics-residual, random-feature,
compressed-input, and classical kernel controls.
\item A closed successor ledger containing exploratory and prospectively
bounded negative candidates, including classical controls that replace the
quantum component while keeping the surrounding construction matched.
\item A self-contained evidence and figure package that supports review of the
negative conclusion without widening the study's scientific claim.
\item A human-centered claim boundary that keeps model evidence, crew-related
requirements, and future human authorization distinct.
\end{enumerate}
